%% kapa: exam=1 round=21 year=2010 period=1 form=A subject=economics subject_ko=경제원론 seq=32 number=72
%% source: Q-Net 감정평가사 기출문제 (cst00309) artlSeq=5024181 file=extracted/1008/5024181/2010년 제21회 감정평가사 1차시험 기출문제/1교시.hwp
%% conversion: kordoc 4.24.51 -> tools/gen1.py (auto); review=proofread-sonnet; flags=none
\begin{question}{72}
적응적 기대가설에 기초한 필립스곡선에 관한 설명으로 옳지 않은 것은?
\begin{choices}
  \item 정부지출이 증가하면 단기적으로 경제의 균형은 필립스곡선을 따라 실업률이 더 낮고 인플레이션율이 더 높은 점으로 옮겨간다.
  \item 통화량이 증가하면 장기적으로 경제의 균형은 필립스곡선을 따라 실업률은 변하지 않고 인플레이션율만 더 높은 점으로 옮겨간다.
  \item 유가상승과 같은 공급충격은 단기적으로 필립스곡선을 왼쪽으로 이동시켜 경제의 균형은 실업률과 인플레이션율이 모두 낮은 점으로 옮겨간다.
  \item 예상인플레이션율이 더 높을수록 단기 필립스곡선은 더 높은 곳에 위치한다.
  \item 프리드만(M. Friedman)에 의하면 장기적으로는 실업률과 인플레이션율 사이에 상충관계가 성립하지 않는다.
\end{choices}
\end{question}
\answer{3}
